\begin{table}[htbp]\centering\small
\caption{Сравнения качества в эксперименте Luna. Разности и границы указаны в процентных пунктах.}
\label{tab:scale-contrasts}
\begin{tabular}{lrrrrr}\toprule
Сравнение & Задач & Разность & 95\%-й интервал & $p$ Холма & Границы\\\midrule
SR -- SN & 982 & $-0.75$ & $[-1.97, 0.44]$ & 0.2357 & $[-2.33, 0.83]$\\
MR -- MN & 964 & $+1.35$ & $[0.17, 2.56]$ & 0.0599 & $[-0.63, 3.37]$\\
MN -- SN & 968 & $-4.17$ & $[-5.65, -2.75]$ & $1.34\times10^{-7}$ & $[-6.00, -2.37]$\\
MR -- SR & 971 & $-2.13$ & $[-3.43, -0.82]$ & 0.0044 & $[-3.83, -0.30]$\\
\bottomrule\end{tabular}
\par\smallskip\noindent Для включения задачи нужны три наблюдаемых повтора в каждом из сравниваемых условий. Границы учитывают все 1\,000 задач и все назначенные повторы; это не доверительные интервалы.
\end{table}
